\documentclass[10pt,a4paper]{article}

% ---------- Packages ----------
\usepackage[margin=0.6in]{geometry}
\usepackage{enumitem}
\usepackage{titlesec}
\usepackage[hidelinks]{hyperref}
\usepackage{amssymb} % for symbol shapes
% \usepackage{fontawesome5}
% \usepackage[T1]{fontenc}

% --- Engine-agnostic font setup ---
\usepackage{ifxetex,ifluatex}
\newif\ifxetexorluatex
\ifxetex \xetexorluatextrue \fi
\ifluatex \xetexorluatextrue \fi

\ifxetexorluatex
  % XeLaTeX / LuaLaTeX: real Gill Sans
  \usepackage{fontspec}
  \defaultfontfeatures{Ligatures=TeX}
  \setmainfont{Gill Sans}[
    UprightFont   = *,
    BoldFont      = {Gill Sans Bold},
    ItalicFont    = {Gill Sans Italic},
    BoldItalicFont= {Gill Sans Bold Italic}
  ]
\else
  % pdfLaTeX fallback: Gill Sans–like
  % Option 1 (closest): Gillius
  \usepackage[sfdefault]{gillius} % needs TeX Live package 'gillius'
  % Option 2: Helvetica
  % \usepackage[scaled=0.95]{helvet}
  \renewcommand{\familydefault}{\sfdefault}
\fi
% ---------- Global Formatting ----------
\pagestyle{empty}
\setlength{\parindent}{0pt}
\setlength{\parskip}{0pt}
\RaggedRight

% Lists: compact
\setlist{nosep,leftmargin=1.2em,labelsep=0.4em}
\setlist[itemize]{topsep=1pt,itemsep=1pt,parsep=0pt}
\renewcommand{\labelitemi}{\tiny$\bullet$}

% Section formatting
\titleformat{\section}{\bfseries\scshape\large}{}{0em}{}[\titlerule]
\titlespacing*{\section}{0pt}{10pt}{5pt}
\titleformat{\subsection}{\bfseries\normalsize}{}{0em}{}
\titlespacing*{\subsection}{0pt}{6pt}{2pt}

% ---------- Header ----------
\newcommand{\cvheader}[6]{%
  \begin{center}
    {\LARGE\bfseries #1}\\[0.9em] 
    \small
    \faEnvelope\ \href{mailto:#2}{#2}\hspace{0.6em}$|$\hspace{0.6em}
    \faGlobe\ \href{https://#3}{#3}\hspace{0.6em}$|$\hspace{0.6em}
    \faGoogle\ \href{#4}{Google Scholar}\hspace{0.6em}$|$\hspace{0.6em}
    % CHANGE IS HERE: Replaced the second {#5} with {GitHub}
    \faGithub\ \href{https://github.com/#5}{GitHub}\hspace{0.6em}$|$\hspace{0.6em}
    {#6}
  \end{center}\vspace{4pt} 
}

% Experience line
\newcommand{\experience}[4]{%
  \vspace{6pt}% extra space before each experience
  \noindent\textbf{#1} \hfill \textit{\small #2}\\
  \textit{#3} \hfill \textit{\small #4}
}

\newcommand{\skill}[2]{%
  \vspace{4pt}
  \textbf{#1: }#2
}

\usepackage{xcolor}
\definecolor{tagpink}{named}{magenta}
\definecolor{taggray}{gray}{0.6} % 0 = black, 1 = white; tweak 0.3 → 0.4/0.5 if you want lighter
\newcommand{\tagRAI}{%
  \textsuperscript{\scriptsize\textcolor{tagpink}{$\blacktriangle$}}%
}
\newcommand{\tagMEAS}{%
  \textsuperscript{\scriptsize\textcolor{tagpink}{$\blacksquare$}}%
}
\newcommand{\tagAGENT}{%
  \textsuperscript{\scriptsize\textcolor{tagpink}{$\bigstar$}}%
}

\begin{document}

\cvheader
{Ruoxi (Anna) Shang}
{rxshang@uw.edu}
{ruoxishang.com}
{https://scholar.google.com/citations?user=}
{ruoxishang}
{Seattle, WA}

% ===== PROFILE (Summary + Research Interests combined) =====
%\section{Research Summary}
%I conduct research at the intersection of HCI and AI. My work focuses on \textbf{evaluating and steering LLM behavior} for trust, alignment, and reliability. I develop theory-grounded \textbf{benchmarks and evaluation pipelines} to measure model risks in realistic settings, and I design human-in-the-loop experiments to audit the downstream impact of AI agents. I have published 10+ peer-reviewed papers (CHI, EMNLP, FAccT, AIES) and contributed to foundational research at Microsoft Research and Oumi AI.


% - understand how and why models succeed and fails as  well as how might we design for better human-ai experiences
% methods:
% - theory-grounded design of framework, strategies, and tasks
% - building pipelines
% - conduct experiments with statistical methods to measure and validate effectiveness
% - I also use HCI methods of studying people as a in-depth way to get subjective yet rich insights on when and why model fails

% ===== EDUCATION =====
\section{Education}
\experience{University of Washington, Human-Centered Design and Engineering}{Sep 2020 -- Exp. Oct 2026}{Ph.D.\ Candidate, Committee: Gary Hsieh, Chirag Shah, Tim Althoff, David Mcdonald}{}
%\\[-2pt]
% \\
% {\small }\vspace{4pt} % clearer advisor line on its own, with spacing

\experience{University of California, Berkeley}{Graduated May 2020}{B.A., Applied Mathematics and Statistics (Concentration: Data Science)}{}

% ===== SKILLS =====

\section{Skills}

\skill{Programming}{Python, JavaScript/TypeScript, React, PyTorch, scikit-learn, SQL, HTML/CSS}

\skill{LLMs}{Prompt Engineering, In-Context Learning, RAG, LLM APIs, Multi-agent frameworks, LLM-as-a-judge, Evaluation and Benchmark Design (e.g. taxonomy development, ablation and robustness test, human annotation)}

\skill{HCI}{Experimental design, Hypothesis testing, Mixed-methods, Survey Design, Think-aloud protocol, Affinity diagramming, Cognitive walkthrough, Heuristic evaluation, Rapid prototyping}



% ===== RESEARCH & WORK EXPERIENCE =====
\section{Research \& Work Experience}

\experience{Oumi AI}{Oct 2025 -- Present}{Research Fellow}{Collaborators: Emmanouil Koukoumidis; Panos Achlioptas; Shang Hong Sim}
\begin{itemize}
  \item Building a \textbf{benchmark for semantic risk reasoning} that evaluates models’ ability to discriminate risk type, assess severity, and justify rationales from text.
  \item Operationalizing a comprehensive risk taxonomy and engineering a high-quality contrastive dataset to stress-test model safety and alignment.
\end{itemize}

% We group your UW work into one "job" but highlight the leadership roles specifically
\experience{University of Washington}{Sep 2020 -- Present}{Research Lead \& Mentor}{Seattle, WA}
\begin{itemize}

  \item Recruited and led a 12-person team to generate complex ground-truth datasets for agent benchmarking. Designed consensus protocols to resolve conflicting domain priorities (Data Science vs. ML) to ensure alignment and rigor. \textcolor{tagpink}{\textit{Co-authored paper accepted to EMNLP 2024}}
  
  \item Advised 50+ interdisciplinary teams through 0-to-1 development of user-centered systems in graduate level project-based classes. Guided the translation of ambiguous user insights into concrete technical specifications, helping teams to evaluate feasibility and manage architectural trade-offs to ensure successful MVP delivery within 10-week constraints.
\end{itemize}

\experience{Microsoft Research}{Jun 2025 -- Sep 2025}{Research Intern}{Advisors: Denae Ford Robinson, Edward Cutrell}
\begin{itemize}
  \item Designed a novel structured \textbf{user profiling framework} that synthesizes interaction data with social science theory to create robust, interpretable user representations, enabling inference-time personalized steering.
  \item Engineered a novel ecologically valid evaluation pipeline to generate ground-truth test cases for auditing agent alignment and validated system effectiveness via a robust auditing protocol using \textbf{90 days of real user communication data} across 20 participants. \textcolor{tagpink}{\textit{Paper and Patent currently under review.}}
\end{itemize}

\experience{Microsoft Research}{Jun 2023 -- Sep 2023}{Research Intern}{Advisors: Steven Drucker; Bongshin Lee; Dave Brown}
\begin{itemize}
  \item Led foundational research on using generative AI for end-to-end data analysis and decision making.
  \item Designed and ran studies with industry experienced data analysts to \textbf{evaluate how they understand, trust, and verify AI-generated analyses}, identifying failure modes and verification strategies to inform the design of AI-assisted data analysis tool. \textcolor{tagpink}{\textit{Co-authored paper accepted to CHI 2024}}
\end{itemize}

\newpage

\experience{TruEra}{Jul 2022 -- Sep 2022}{Data Science and User Research Intern}{Mentors: Mantas Lilis; Joshua Noble; Justin Lawyer}
\begin{itemize}
  \item Improved the Machine Learning (ML) model (e.g. classification, regression) debugging experience for data scientists by leading research and designing product guided workflows from research insights for TruEra \textbf{ML Model Diagnostics}. Identified the top pain points in real debugging practice and prioritized fixes, driving 0 to 1 product exploration during after new funding round and shaped the roadmap.
\end{itemize}


% ===== PUBLICATIONS =====
\section{Selected Publications}
{\small
\textit{
$\textcolor{tagpink}{\blacktriangle}$ Responsible AI \quad
$\textcolor{tagpink}{\blacksquare}$ Measurement \& Benchmarks \quad
$\textcolor{tagpink}{\bigstar}$ Agents \& LM-powered Systems}
}

\vspace{6pt}


\begin{itemize}[label={},leftmargin=0em,itemindent=0em,parsep=1pt]

  \item \textbf{Ruoxi Shang}, Dan Marshall, Edward Cutrell, Denae Ford Robinson.
  \textit{Mimetic Alignment with PRISMA: Auditing of AI-inferred Personal Profiles}.
  \textit{arXiv} (under review).\hfill\tagRAI\tagAGENT
  \vspace{8pt}
  
  \item Meziah Ruby Cristobal*, Hyeonjeong Byeon*, Tze-Yu Chen*, \textbf{Ruoxi Shang}*, Donghoon Shin*, Ruican Zhong*, Tony Zhou*, Gary Hsieh.
  \textit{PaperTok: Exploring the Use of Generative AI for Creating Short-form Videos for Research Communication}.
  \textit{arXiv} (under review). {\small[* equal contribution]}\hfill\tagAGENT
  \vspace{8pt}
  
  \item \textbf{Ruoxi Shang}, Ruican Zhong, Yuren Pang, David McDonald, Chirag Shah, Gary Hsieh.
  \textit{Synthetic Diversity: How Researchers Perceive and Engage with LLM-Generated Diverse Research Feedback}.
  \textit{arXiv} (under review).\hfill\tagRAI\tagMEAS
  \vspace{8pt}
  
  \item \textbf{Ruoxi Shang}, Lei Cai, Courtney Cho, Chirag Shah, Gary Hsieh.
  \textit{Creative Stimuli for Multi-Agent LLM Brainstorming}.
  \textit{arXiv} (under review).\hfill\tagAGENT
  \vspace{8pt}

  \item \textbf{Ruoxi Shang}, Keri Mallari, Wei Bin Au Yeong, Ken Yasuhara, Anthony Tang, Gary Hsieh.
  \textit{Rethinking Teaching Evaluation Reports}.
  \textit{CSCW 2025}.\hfill\tagRAI\tagMEAS
  \vspace{8pt}

  \item \textbf{Ruoxi Shang}, Gary Hsieh, Chirag Shah.
  \textit{Trusting your AI Agent Emotionally and Cognitively: Development and Validation of a Semantic Differential Scale for AI Trust}.
  \textit{AIES 2024}.\hfill\tagRAI\tagMEAS
  \vspace{8pt}

  \item Ken Gu, \textbf{Ruoxi Shang}, Ruien Jiang, Keying Kuang, \emph{et al.}
  \textit{BLADE: Benchmarking Language Model Agents for Data-Driven Science}.
  \textit{EMNLP 2024}.\hfill\tagMEAS\tagAGENT
  \vspace{8pt}

  \item Ken Gu, \textbf{Ruoxi Shang}, Tim Althoff, Chenglong Wang, Steven M.\ Drucker.
  \textit{How Do Analysts Understand and Verify AI-Assisted Data Analyses?}
  \textit{CHI 2024}.\hfill\tagRAI\tagAGENT
  \vspace{8pt}

  \item \textbf{Ruoxi Shang}, Kevin Feng, Chirag Shah.
  \textit{Lay Users' Needs of Counterfactual Explanations}.
  \textit{FAccT 2022}.\hfill\tagRAI
  \vspace{8pt}

\end{itemize}

\end{document}